\section{Workstream 2: processing and feature engineering}
\label{sec:processing}

\subsection{DS1: from raw pages to a clean transaction table}

\code{python -m explorer.cli transactions} turns the raw pages into
\fp{data/processed/esplora/transactions.csv}, one row per (address, transaction). Each step logs
how many rows it removed, and that log becomes the processing table in the report:

\begin{enumerate}
  \item Parse every page and keep confirmed transactions only. The chain endpoint returns only
    confirmed ones, but the check is explicit anyway.
  \item Remove duplicate (address, txid) pairs, which can occur at page boundaries if a new block
    arrives mid-fetch.
  \item Derive \code{received_sats}, \code{sent_sats}, \code{net_sats} and \code{direction} with
    \code{amounts.address_flow}. Also derive \code{n_inputs}, \code{n_outputs}, \code{fee_sats},
    \code{vsize_vb} and \code{block_time_utc}.
  \item Drop transactions where the address neither receives nor sends (expected count: 0).
    A non-zero count is logged as a data problem.
  \item Sort by (address, block time, txid). Transactions in the same block share a time and are
    never part of each other's history.
  \item Units are always in the column names (\code{_sats}, \code{_utc}, \code{_hours}). Amounts
    stay integers in satoshis; BTC appears only in plots and in the interface.
\end{enumerate}

\textbf{Missing values}, recorded per column. The gap (F3) doesn't exist for an address's first
transaction, so it stays empty rather than being filled with an invented number, and the
transaction isn't scored. Fees are always present for confirmed transactions. Any other missing
value stops the pipeline with a clear message, because it would mean the API changed.

\subsection{Features F1--F4: a port of the prototype}
\label{sec:features}

\fp{explorer/features.py} is a line-by-line Python port of \code{scoring.ts} and reads the same
\code{settings.json}. That keeps A1 §2.3, the prototype and the pipeline consistent.

\begin{center}
\small
\begin{tabularx}{\textwidth}{@{}L{1.0cm}L{2.1cm}Y L{3.3cm}@{}}
\toprule
\textbf{ID} & \textbf{Feature} & \textbf{Definition (history = strictly earlier transactions of the same address)} & \textbf{Scale for the z-score} \\
\midrule
F1 & Size & $\log_{10}(1 + |\text{net}|)$, in satoshis & as is \\
F2 & Burst & number of earlier transactions in the previous 7 days & as is (count) \\
F3 & Gap & hours since the previous earlier transaction & $\log_{10}(1+h)$ \\
F4 & Relative size & $|\text{net}|$ ÷ median $|\text{net}|$ of the size window & $\log_{10}(\max(r, 0.001))$ \\
\bottomrule
\end{tabularx}
\end{center}

\textbf{Comparison windows} (CR-02): the previous 30 days, extended back until at least 20 earlier
transactions are included. The \emph{timing window} holds all earlier transactions and is used for
F2 and F3. The \emph{size window} holds only earlier transactions in the same direction (CR-04)
and is used for F1 and F4.

\textbf{Robust z} for each feature, against the same feature over its window, with a floor on the
spread (CR-05):
\[
z = \frac{x - \operatorname{median}(W)}{\max\bigl(1.4826 \cdot \operatorname{MAD}(W),\ s_{\min}\bigr)},
\qquad s_{\min} = (0.1,\ 1,\ 0.1,\ 0.1)\ \text{for F1--F4}.
\]

\textbf{Scorable} means at least 10 earlier transactions spanning at least 14 days, and at least
5 values in the window. Unscorable transactions are kept and displayed, but never flagged.

\textbf{Extended features} (used only in an ablation and as TxGPT tokens, never in the
Explorer's explanation): F5 $\log_2 n_{\text{inputs}}$, F6 $\log_2 n_{\text{outputs}}$, F7
$\log_{10}$ fee rate (sat/vB), and F8 the hour of day as (sin, cos).

\textbf{Inputs to pooled models.} One model is fitted across all 25 addresses, which hold very
different amounts: a tip jar next to a merchant. So the clustering, IF and LOF inputs are each
transaction's \emph{own-history z-scores} $(z_1,\dots,z_4)$, clipped to ±20 so that a single
extreme value can't dominate, plus a one-hot direction. This makes ``unusual for this address''
comparable across addresses. The ablation repeats the models on transformed raw features scaled
with a \code{StandardScaler} fitted on the fit period.

\textbf{Refinement ablation (CR-02, CR-04, CR-05).} The baseline and every downstream model run
twice: once with the method as A1 §2.3 wrote it (a 30-day window, all directions, no transform,
no floor), and once with the refinements. For each run the evidence recorded is the injection
recall per pattern, the share of addresses that break the 20\% sanity rule, and the flag rate.
The team adopts the refinements in the Change Log only if the evidence supports them. Either way,
the report states which version is final and why.

\subsection{The time split}
\label{sec:split}

For each address, take the scorable transactions in time order. The cut-off time $t^\ast$ is the
time of the transaction at the 70th percentile (\code{fitFraction} = 0.7). Transactions with time
$< t^\ast$ form the \textbf{fit period}; those with time $\geq t^\ast$ form the \textbf{evaluation
period}. Same-block transactions therefore never straddle the cut. Unscorable early transactions
serve only as history. Every model is fitted on fit-period rows pooled across addresses. Every
threshold, such as the top-2\% cut, is set from fit-period scores and then applied unchanged to
the evaluation period.

\subsection{Synthetic injection: known positives in real histories}
\label{sec:injection}

\code{explorer/inject.py} plants patterns in a \emph{copy} of one address's history, inside its
evaluation period. Features are then recomputed time-safely, because a planted transaction
changes the features of everything after it. Finally the copy is scored with \emph{frozen}
models. A single plant goes into each copy, so plants never mask one another.

\begin{center}
\small
\begin{tabularx}{\textwidth}{@{}L{2.4cm}Y@{}}
\toprule
\textbf{Pattern (A1 §2.3.2)} & \textbf{How it is planted} \\
\midrule
Value spike & One transaction at a random time between two real ones. Its direction is the address's majority direction over the previous 30 days. Its amount is 10 × that direction's median $|\text{net}|$ up to that point. Input and output counts are the address's median. \\
Burst & Ten transactions within a random 24-hour period, with amounts drawn from the address's own earlier amounts, so only the timing is unusual. Scored two ways: per transaction (A1's measure, the share of the 10 flagged) and per event (at least one flagged). \\
Dormant wake-up & One transaction appended 60--120 days (uniform) after the address's last real one, with a typical amount, so only the gap is unusual. \\
\bottomrule
\end{tabularx}
\end{center}

There are 50 plants per pattern, spread across addresses in proportion to their evaluation-period
size, and the whole set is repeated with 5 seeds (30049 to 30053). Results are reported as mean ±
standard deviation. Each planted row carries \code{is_synthetic = 1} and can never leak into a
fit set.

\subsection{DS3 Elliptic processing}

\begin{enumerate}
  \item Join features to classes on \code{txId} (a left join; every feature row has a class row).
    Map \code{1} to illicit (1), \code{2} to licit (0), and \code{unknown} to missing.
  \item Name the columns \code{time_step}, \code{local_01}--\code{local_93} and
    \code{agg_01}--\code{agg_72}.
  \item \textbf{Join the edge list} to add two engineered graph features, in-degree and
    out-degree per transaction. This is a genuine dataset join, one of the processing steps the
    brief lists.
  \item Feature sets: \textbf{LF}, the 93 local features (+2 degrees), and \textbf{AF}, all 165
    (+2). \code{time_step} is never a feature; it only defines the split.
  \item \textbf{Temporal split} as in Weber et al.\ (2019): time steps 1--34 for training and
    35--49 for testing. Tuning uses only training time steps: fit on 1--26, validate on 27--34.
  \item Scaling inside a scikit-learn \code{Pipeline}, so the scaler only ever sees training
    rows. Check for missing and constant columns and drop constants.
\end{enumerate}

\subsection{DS4 BitcoinHeist processing}

\begin{enumerate}
  \item Binary label: \code{white} → 0, any family name → 1. The family is kept for the
    leave-family-out test.
  \item $\log_{10}$ \code{income} (satoshis). Apply $\log(1+x)$ to \code{weight}, \code{count},
    \code{looped} and \code{neighbors} if the analysis confirms heavy tails.
  \item \code{address}, \code{year} and \code{day} are \emph{not} features. The address is an
    identifier. Year and day would let a model learn when each family was active instead of what
    the activity looked like, which is a leakage shortcut. They only define splits.
  \item Remove exact duplicate rows and record how many were removed.
  \item \textbf{Splits}, fixed after the analysis by a written rule. (a) Temporal: train on years
    up to $Y$ and test on later years, with $Y$ the latest year that leaves at least 1,000
    positives in the test set. (b) A stratified random 80/20 split as an optimistic reference.
    (c) Leave one family out, for each of the five largest families.
  \item \textbf{Training-set sampling.} All positives plus negatives downsampled to a 10:1 ratio
    (seeded), \emph{in the training set only}. Test sets keep the real class ratio, because
    precision depends on the base rate.
\end{enumerate}

\subsection{DS2: TxGPT tokenisation}
\label{sec:txtokens}

Each transaction becomes 6 tokens, always in this order. Gap comes first so the model predicts
the transaction's content knowing when it happened:

\begin{center}
\small
\begin{tabularx}{\textwidth}{@{}L{1.5cm}C{1.2cm}Y@{}}
\toprule
\textbf{Field} & \textbf{Tokens} & \textbf{Buckets} \\
\midrule
GAP & 18 & FIRST (no earlier transaction); SAME\_BLOCK; then hours since the previous transaction, with upper bounds 1/6, 1, 3, 6, 12, 24, 48, 96, 168, 336, 720, 1,440, 2,160, 4,320, 8,760, and above \\
DIR & 3 & IN, OUT, ZERO (net 0) \\
AMT & 47 & $\lfloor \log_2 |\text{net}| \rfloor$ clipped to 0--46: one bucket per doubling, from 1 sat to about 700,000 BTC \\
NIN & 7 & input count: 1, 2, 3--4, 5--9, 10--19, 20--49, 50 or more \\
NOUT & 7 & output count, same buckets \\
FEE & 9 & fee rate in sat/vB: below 1, then doubling bounds 2, 4, 8, 16, 32, 64, 128, and above \\
Special & 3 & PAD, BOS (start of a history), UNK \\
\midrule
Total & 94 & padded to 128 for GPU-friendly shapes (\code{pad_vocab} in \code{lm/model.py}) \\
\bottomrule
\end{tabularx}
\end{center}

\begin{itemize}
  \item A history becomes BOS followed by $6n$ tokens, oldest first.
  \item Training windows are 768 tokens (128 transactions). Longer histories are cut into windows
    with a stride of 384 tokens, aligned to transaction boundaries. BOS appears only at a true
    history start, so the model learns both ``beginning'' and ``middle'' contexts.
  \item Split \emph{by address}: 90\% train, 5\% validation, 5\% test. DS1 is never in DS2.
  \item Written by \code{python -m lm.prepare_tx} as \code{train.bin}, \code{val.bin} and
    \code{test.bin} (uint8 is enough for 128 tokens), plus \code{meta.json} with the bucket edges,
    so scoring uses identical buckets.
\end{itemize}

\subsection{Leakage controls in one place}
\label{sec:leakage}

\begin{enumerate}
  \item Features use strictly earlier transactions only, and same-block transactions are
    excluded (DR-02).
  \item The fit/evaluation split is by time within each address. Models and thresholds are fitted
    on the fit period only.
  \item Every scaler lives inside a \code{Pipeline} fitted on training rows.
  \item Elliptic uses a temporal split, with tuning confined to training time steps. The test
    period is used once.
  \item BitcoinHeist uses neither year nor day as a feature; downsampling happens in the training
    set only.
  \item TxGPT and DS7 are split by address, and DS1 is excluded from DS2.
  \item Plants go only into evaluation-period copies, scored by frozen models.
  \item No test number is ever used to choose a hyperparameter. If one is, the report says so.
\end{enumerate}

\subsection{Tests}
\label{sec:tests}

\begin{center}
\small
\begin{tabularx}{\textwidth}{@{}L{4.6cm}Y@{}}
\toprule
\textbf{Test file} & \textbf{What it proves} \\
\midrule
\code{test_amounts.py} (existing, 9) & Received, sent, net and direction on hand-made transactions. \\
\fp{test_features_time_safe.py} & \textbf{DR-02 condition 1}: for 20 random transactions, features recomputed from a history truncated to strictly earlier transactions equal the pipeline's values exactly. \\
\code{test_features_parity.py} & The Python port matches \code{scoring.ts} on the prototype's illustrative histories, to within $10^{-9}$. A small script dumps the TypeScript results to JSON for this. \\
\code{test_split.py} & \textbf{DR-02 condition 2}: no txid is in both periods, and the latest fit time is before the earliest evaluation time for every address. \\
\code{test_injection.py} & Plants land in the evaluation period only. A spike is 10 × the median, a burst is 10 transactions within 24 hours, and a wake-up gap is at least 60 days. Every plant is marked synthetic. \\
\code{test_tx_tokens.py} & Bucket edges, token round trips, a sequence length of $1+6n$, and an identical vocabulary at training and scoring time. \\
\code{test_predict.py} & Prediction outputs match the A1 API shapes, the limitation line is present on every flagged item, and no banned word appears in any user-facing string. \\
\code{tests/test_lm.py} (existing, 15) and \code{lm.smoke} & The language-model code end to end on CPU. \\
\bottomrule
\end{tabularx}
\end{center}

All of these run with \code{python -m pytest -q} in under two minutes on CPU, which is also the
README's quick check.

\begin{team}
A1 requires DR-02 code review ``by a member other than the author''. Le Quang Anh Hoang reviews
\code{features.py}, \code{split.py} and \code{inject.py} on 24--25 October, against the tests
above, and their review comments go in the minutes.
\end{team}

\begin{feeds}
Report §4 (Data processing): the cleaning log table, the feature table, the time split, the
leakage list (condensed to 3--4 sentences), the DS3 join, the DS4 sampling and the TxGPT
tokenisation. Code rubric lines 1, 4 and 5, through testable, documented modules.
\end{feeds}
